\subsection{Majorizable vectorial measures}

Let \(q\) be a lower semi-continuous semi-norm on \(F\) . We shall denote by \(A_q'\) the set of \(z' \in F'\) such that \(|\langle z', x \rangle| \leqslant q(x)\) for all \(x \in F\) . This is the polar in \(\mathbf{F}'\) of the set of \(\mathbf{x} \in \mathbf{F}\) such that \(q(\mathbf{x}) \leqslant 1\) ; for every \(\mathbf{x} \in \mathbf{F}\) , \(q(\mathbf{x}) = \sup_{\mathbf{z}' \in \hat{\mathrm{A}}_q'} |\langle \mathbf{x}, \mathbf{z}' \rangle|\) .

DEFINITION 3. --- Let m be a vectorial measure on T with values in F. If q is a lower semi-continuous semi-norm on F, m is said to be q-majorizable if there exists a positive measure \(\mu\) such that \(|z^{\prime} \circ m| \leqslant \mu\) for every \(z^{\prime} \in A_{q}^{\prime}\) ; the supremum of the measures \(|z^{\prime} \circ m|\) as \(z^{\prime}\) runs over \(A_{q}^{\prime}\) (Ch. III, §2, No.~4, Th. 3) is then denoted \(q(\mathbf{m})\) . One says that m is majorizable if it is q-majorizable for every continuous semi-norm q on F.

If \(\mathbf{m}\) and \(\mathbf{m}'\) are both \(q\) -majorizable, it is immediate that \(\mathbf{m} + \mathbf{m}'\) is also \(q\) -majorizable and that

\[
q (\mathbf {m} + \mathbf {m} ^ {\prime}) \leqslant q (\mathbf {m}) + q (\mathbf {m} ^ {\prime}).
\]

When F is a normed space, with norm denoted \(|x|\) , to say that m is majorizable therefore means that the measures \(|z'\circ m|\) , where \(|z'| \leqslant 1\) , are bounded above \(^{1}\) by a same positive measure; one then denotes by \(|m|\) the supremum of this family of measures.

If \(\mathbf{F} = \mathbf{R}\) , the measure \(|m|\) corresponding to the Euclidean norm \(|x|\) on \(\mathbf{R}\) coincides with the measure denoted by \(|m|\) in Ch. III, §1, No.~5.

PROPOSITION 4. --- Let \((\mathrm{F}_i)_{1 \leqslant i \leqslant n}\) be a finite family of Hausdorff locally convex spaces, \(\mathbf{F} = \prod_{i=1}^{n} \mathbf{F}_i\) their product, \(q_i\) ( \(1 \leqslant i \leqslant n\) ) a lower semi-continuous semi-norm on \(\mathbf{F}_i\) , and \(q\) the semi-norm on \(\mathbf{F}\) defined by \(q(\mathbf{x}_i, \ldots, \mathbf{x}_n) = \sum_{i=1}^{n} q_i(\mathbf{x}_i)\) . If \(\mathbf{m}_i\) ( \(1 \leqslant i \leqslant n\) ) is a vectorial measure on \(\mathbf{T}\) with values in \(\mathbf{F}_i\) that is \(q_i\) -majorizable, then the measure \(\mathbf{m} = (\mathbf{m}_1, \ldots, \mathbf{m}_n)\) with values in \(\mathbf{F}\) is \(q\) -majorizable.

For, the dual \(\mathbf{F}'\) may be identified with \(\prod_{i=1}^{n} \mathrm{F}_i'\) , in such a way that if \(\mathbf{x} = (\mathbf{x}_i) \in \mathrm{F}\) , \(\mathbf{z}' = (\mathbf{z}_i') \in \mathrm{F}'\) , then \(\langle \mathbf{x}, \mathbf{z}' \rangle = \sum_{i=1}^{n} \langle \mathbf{x}_i, \mathbf{z}_i' \rangle\) . If \(|\langle \mathbf{x}, \mathbf{z}' \rangle| \leqslant q(\mathbf{x})\) for every \(\mathbf{x} \in \mathrm{F}\) , then in particular \(|\langle \mathbf{x}_i, \mathbf{z}_i' \rangle| \leqslant q_i(\mathbf{x}_i)\) for \(1 \leqslant i \leqslant n\) , and the converse is obvious, which shows that the set \(\mathrm{A}_q'\) is the product of the \(\mathrm{A}_{q_i}'\) . Since by hypothesis \(|\mathbf{z}_i' \circ \mathbf{m}_i| \leqslant q_i(\mathbf{m}_i)\) for \(\mathbf{z}_i' \in \mathrm{A}_{q_i}'\) , it follows that

\[
\left| \mathbf {z} ^ {\prime} \circ \mathbf {m} \right| \leqslant \sum_ {i = 1} ^ {n} \left| \mathbf {z} _ {i} ^ {\prime} \circ \mathbf {m} _ {i} \right| \leqslant \sum_ {i = 1} ^ {n} q _ {i} (\mathbf {m} _ {i})
\]

for every \(\mathbf{z}' \in \mathrm{A}_q'\) , which proves the proposition.

COROLLARY. --- If the space F is finite-dimensional, then every vectorial measure m with values in F is majorizable. In order that a numerical function be essentially integrable for m, it is necessary and sufficient that it be essentially integrable for \(|m|\) (where \(|x|\) denotes any norm on F).

PROPOSITION 5. --- Let q be a lower semi-continuous semi-norm on F. Let m be a q-majorizable measure, and let f be a function essentially integrable for m and such that \(\int f dm \in F\) . Then

\[
q \left(\int f d \mathbf {m}\right) \leqslant \int^ {\bullet} | f | d q (\mathbf {m}).
\]

For,

\[
q \left(\int f d \mathbf {m}\right) = \sup _ {\mathbf {z} ^ {\prime} \in \mathrm{A} _ {q} ^ {\prime}} \left| \left\langle \mathbf {z} ^ {\prime}, \int f d \mathbf {m} \right\rangle \right| \leqslant \sup _ {\mathbf {z} ^ {\prime} \in \mathrm{A} _ {q} ^ {\prime}} \int | f | d | \mathbf {z} ^ {\prime} \circ m | \leqslant \int^ {\bullet} | f | d q (\mathbf {m})
\]

by virtue of (1) and the relation \(|\mathbf{z}' \circ \mathbf{m}| \leqslant q(\mathbf{m})\) for \(\mathbf{z}' \in \mathrm{A}_q'\) .

PROPOSITION 6. --- Let F be a quasi-complete, Hausdorff locally convex space, m a majorizable measure on T with values in F. If f is a numerical function essentially integrable for all of the measures \(q(\mathbf{m})\) (where q runs over the set of continuous semi-norms on F), then f is essentially integrable for m, and \(\int f dm \in F\) .

We shall make use of the following auxiliary result. Let \((\mu_{\iota})_{\iota\in\mathrm{I}}\) be an increasing directed family of positive measures on T. Let us denote by \(\mathcal{L}^{1}\big((\mu_{\iota})_{\iota\in\mathrm{I}}\big)\) the vector space of finite numerical functions on T, essentially \(\mu_{\iota}\) -integrable for every \(\iota\in I\) , equipped with the topology defined by the semi-norms \(f\mapsto\mu_{\iota}(|f|)\;\ (\iota\in\mathrm{I})\) . Let \(L_{0}\) be the linear subspace of \(\mathcal{L}^{1}\big((\mu_{\iota})_{\iota\in\mathrm{I}}\big)\) generated by the products \(g\varphi_{K}\) , where g runs over the set of continuous finite numerical functions on T, and K over the set of compact subsets of T.

Lemma 1. --- When \(\mathcal{L}_0\) and \(\mathcal{K}(\mathrm{T})\) are equipped with the topology induced by that of \(\mathcal{L}^1\big((\mu_\iota)_{\iota \in \mathrm{I}}\big)\) :

\begin{enumerate}
\def\labelenumi{\alph{enumi})}
\item
  each element of \(\mathcal{L}_0\) is in the closure of some bounded subset of \(\mathcal{K}(\mathrm{T})\) ;
\item
  each element of \(\mathcal{L}^1\big((\mu_\iota)_{\iota \in \mathrm{I}}\big)\) is in the closure of some bounded subset of \(\mathcal{L}_0\) .
\end{enumerate}

To prove a), we may restrict ourselves to the case of an element of the form \(f = g\varphi_{\mathrm{K}}\) ( \(g \in \mathcal{C}(\mathrm{T})\) ), K compact in T). It is immediate (by virtue of Urysohn's theorem) that \(f\) is in the closure of the set B of functions of the form \(gh\) , where \(h\) describes the set of continuous mappings of T into {[}0, 1{]} that are equal to 1 on K and to 0 on the complement of a fixed compact neighborhood H of K. Moreover, the set B is bounded, because \(\mu_{\iota}(|gh|) \leqslant \mu_{\iota}(|f\varphi_{\mathrm{H}}|)\) for every function h having the preceding properties.

Let us now prove b); we may restrict ourselves to the case of an element \(f \geqslant 0\) of \(\mathcal{L}^1((\mu_\iota)_{\iota \in \mathrm{I}})\) . For every \(\iota \in \mathrm{I}\) and every \(\varepsilon > 0\) , there exists a compact subset \(\mathrm{K}(\iota, \varepsilon)\) of \(\mathrm{T}\) such that the restriction of \(f\) to \(\mathrm{K}(\iota, \varepsilon)\) is continuous and \(|\mu_\iota(|f - f\varphi_{\mathrm{K}(\iota, \varepsilon)}|) \leqslant \varepsilon\) . It is clear that \(f\) is in the closure of the set \(\mathrm{C}\) of \(f\varphi_{\mathrm{K}(\iota, \varepsilon)}\) (where \(\iota \in \mathrm{I}, \varepsilon > 0\) ). By virtue of Urysohn's theorem, the set \(\mathrm{C}\) in contained in \(\mathcal{L}_0\) ; moreover, it is bounded, because \(\mu_\varkappa(f\varphi_{\mathrm{K}(\iota, \varepsilon)}) \leqslant \mu_\varkappa(f)\) for all \(\iota \in \mathrm{I}, \varkappa \in \mathrm{I}\) and \(\varepsilon > 0\) .

Let us now prove Prop. 6: for every function \(g \in \mathcal{K}(\mathrm{T})\) and every continuous seminorm \(q\) on \(\mathrm{F}\) , \(q\big(\int g dm\big) \leqslant \int |g| d\big(q(\mathbf{m})\big)\) (Prop. 5), which implies that the mapping \(g \mapsto \int g dm\) of \(\mathcal{K}(\mathrm{T})\) into \(\mathrm{F}\) is continuous when \(\mathcal{K}(\mathrm{T})\) is equipped with the topology induced by that of \(\mathcal{L}^1\big((q(\mathbf{m}))_{q \in \mathbb{Q}}\big)\) (Q the set of continuous semi-norms on \(\mathrm{F}\) ). Consequently, by the preceding lemma and Prop. 10 of TVS, III, §1, No.~6, this mapping may be extended by continuity, first to a continuous linear mapping \(v_0\) of \(\mathcal{L}_0\) into \(\mathrm{F}\) , then to a continuous linear mapping \(v\) of \(\mathcal{L}^1\big((q(\mathbf{m}))_{q \in \mathbb{Q}}\big)\) into \(\mathrm{F}\) . Moreover, for every \(\mathbf{z}' \in \mathrm{F}'\) the relation \(\langle \mathbf{z}', v(f) \rangle = \int f d(\mathbf{z}' \circ \mathbf{m})\) holds, by the definition of \(v\) , for every \(f \in \mathcal{K}(\mathrm{T})\) ; since \(|\mathbf{z}' \circ \mathbf{m}| \leqslant q(\mathbf{m})\) for \(q(\mathbf{z}) = |\langle \mathbf{z}', \mathbf{z} \rangle|\) , the mapping \(f \mapsto \int f d(\mathbf{z}' \circ \mathbf{m})\) is continuous on \(\mathcal{L}^1\big((q(\mathbf{m}))_{q \in \mathbb{Q}}\big)\) , therefore again by continuity, the relation \(\langle \mathbf{z}', v(f) = \int f d(\mathbf{z}' \circ \mathbf{m})\) holds for every function \(f \in \mathcal{L}^1\big((q(\mathbf{m}))_{q \in \mathbb{Q}}\big)\) . It follows that \(\nu(f) = \int f dm\) , which completes the proof.

